\section{Definitions and premises}\label{sec:classification}

\medskip\noindent\textbf{D1. Worldview.}\phantomsection\label{item:D1}
(Definition.) A claim about what ultimately is. Every worldview goes beyond any finite body of
evidence. Examples are theism, naturalism, the skeptical hypothesis (that we are deceived, or
simulated), and the teaching of any religious tradition.

\medskip\noindent\textbf{D2. Classic argument.}\phantomsection\label{item:D2}
(Definition.) An argument offered to establish, or to refute, a worldview.

\medskip\noindent\textbf{D3. Establish.}\phantomsection\label{item:D3}
(Definition.) An argument establishes a claim from the evidence if and only if the claim is true in every
account consistent with the evidence and with the truths of reason. This is \emph{Algorithmic Agnosticism}'s
standard for a theory known to be final, applied to worldviews: the claim holds in every account the evidence leaves open.

Natural theology claims demonstration. Aquinas holds that the existence of God, though not
self-evident to us, can be demonstrated from its effects, which are known to us
\citep[I q.2 a.2]{aquinas-st}. The First Vatican Council teaches that God, the beginning and end of
all things, ``can be known with certainty'' from created things by the natural light of human
reason \citep[ch.~2]{vatican1870}. The denial has had its demonstrations too. The logical problem
of evil holds that God's existence is inconsistent with the existence of evil, and Findlay argued
that the being worship requires would have to exist necessarily, and that, if necessity belongs
only to propositions, no being can \citep{mackie1955,findlay1948}. Each side claims more than
probability. D3 names the standard that these proofs and disproofs alike set
out to meet.

\medskip\noindent\textbf{P-M. Mapping premise.}\phantomsection\label{item:PM}
(Premise.)
\begin{enumerate}[label=(\alph*)]
  \item A worldview, insofar as it claims to be true, is an account in the sense of
        \emph{Algorithmic Agnosticism}. Every route to belief, argument or not, supplies evidence, a judgment, or an
        imposed constraint.
  \item For any finite evidence, any account that goes beyond it, a worldview included, has a
        rival consistent with the evidence and the truths of reason on which it is false.
\end{enumerate}
\emph{Motivation.} Clause (a) follows the treatment of religious belief as a ``religious
hypothesis'' \citep[§X]{james1896}, or as an explanatory hypothesis assessed by its fit with the
evidence \citep{swinburne2004}. \emph{Algorithmic Agnosticism} proves for its own accounts what (b) asserts of
worldviews (AA-R1), with rivals that differ on what has not yet been observed. Those accounts do
not speak of God or nature, so for worldviews (b) is a premise.
Two rivals illustrate (b). For theism, one rival is an account on which the same course of
experience, religious experience included, arises with nothing beyond nature: the sense of God's
presence, the order of the world and the testimony of the traditions are what natural processes
would bring about. For naturalism, one rival is an account on which God creates and sustains
exactly the world the sciences describe: every law holds and every experience runs as it does,
but the world depends for its existence on God. Each rival agrees with all the evidence that the
view it opposes cites. Each is the kind of rival the classic debates already treat as live: theists answer
natural explanations of religious experience by argument, and naturalists answer a God who works
through the order of nature by argument, rather than dismissing either as incoherent. What (b)
asserts of every worldview is thus what each side already treats as live in the other's: that it
has such a rival.
\emph{Refuted by:} for (a), a worldview that claims to be true yet is no account of what is, or a
route to belief that delivers something else; for (b), an account that goes beyond some finite evidence, a worldview included, yet is
true in every account consistent with that evidence and the truths of reason. A sound ontological
argument would show this for theism, and a sound logical problem of evil for its denial.

\medskip\noindent\textbf{T-Class. Classification premise.}\phantomsection\label{item:TClass}
(Premise; used for coverage in \cref{sec:arguments}.) Every element of a classic argument is one
of five things: evidence, an account, a judgment, an imposed constraint, or a truth of reason.
\begin{itemize}
  \item \emph{Propositions} \citep[§4]{hume1748}.
        A matter of fact is either a report of experience, which is evidence, or a claim about
        what is, which is an account. A relation of ideas is a truth of reason. A claim offered
        \emph{a priori} that would exclude an account consistent with the truths of reason is not
        a truth of reason, since that account is consistent with them. It serves as a criterion for
        choosing among accounts the evidence leaves open, so it is an imposed constraint.
        Examples are the possibility of a maximally great being, the principle of
        sufficient reason, the incompatibility of God and evil, and a disputed modal axiom.
  \item \emph{Rules.} A valid deductive rule is a truth of reason. A non-deductive rule
        (induction, inference to the best explanation, probabilistic confirmation, the
        deliverance of a faculty) is a judgment. A standard for preferring one conclusion to
        another is an imposed constraint.
  \item \emph{Practical reasons}, such as those of \citet{pascal1670} and \citet{james1896},
        are reasons for choosing a belief, so they are imposed constraints.
  \item \emph{Nothing else enters an argument.} An argument consists of propositions, and of
        rules by which they are taken to support a conclusion.
\end{itemize}
\emph{Refuted by} an element of a classic argument that fits none of the five classes. Hume's fork
is itself contested, so T-Class is owned as a premise, not claimed as a truth of reason.
